% sec_verified_A2_results.tex -- independently verified structural extensions
\section{Mode Closure, Gaussian Compression, and Dimensional Thresholds}
\label{sec:verified-A2-results}

This section records extensions which do not follow by merely evaluating
further members of the coefficient hierarchy.  The comparison law is allowed
to move, and the multidimensional argument is carried out in the topology
actually needed for the leading entropy coefficient.  Throughout this
section, Gaussian quadrature refers to quadrature for a probability measure
on the real line, not to Gaussian smoothing.

\subsection{The first Cayley channel}

Let \(X=\gamma-\bar\gamma\), and, for \(t>0\), define
\begin{equation}\label{eq:first-cayley-channel-definition}
 a_t(X):=\frac{\ii X}{2t-\ii X},
 \qquad c_k(t):=\EE[a_t(X)^k],\quad k\ge1.
\end{equation}
The following statement concerns the centred law of \(X\).  In particular,
it is invariant under a simultaneous translation of the original law and
its matched reference.

\begin{theorem}[First-channel closure and the even-moment hierarchy]
\label{thm:first-cayley-mode-closure}
Let \(\nu\) be a probability law on \(\RR\) with finite mean, and use the
Cayley coordinate
\[
 z=\frac{1+\ii y}{1-\ii y}\in\TT\setminus\{-1\}.
\]
Then, for every \(t>0\),
\begin{equation}\label{eq:first-channel-exact-fourier-series}
 \frac{(P_t*\nu)(\bar\gamma+ty)}{P_t(ty)}
 =1+\sum_{k\ge1}
 \bigl(\overline{c_k(t)}z^k+c_k(t)z^{-k}\bigr).
\end{equation}
The series converges pointwise for every \(y\in\RR\), locally uniformly in
\(y\), and in \(L^2(\omega)\) whenever \(\EE X^2<\infty\).  Its coefficients
satisfy the exact differential closure
\begin{equation}\label{eq:first-channel-differential-closure}
 c_{k+1}(t)=-c_k(t)-\frac{t}{k}c_k'(t),
 \qquad k\ge1.
\end{equation}
Consequently, the scale trajectory \(c_1\), with its successive derivatives,
determines the matched quotient of the centred law at every scale.  Under
finite variance one also has
\begin{equation}\label{eq:first-channel-chi-square-parseval}
 \chi^2\!\left(P_t*\nu\,\middle\|\,
 P_t(\mathord\cdot-\bar\gamma)\right)
 =2\sum_{k\ge1}|c_k(t)|^2.
\end{equation}
If \(\varphi_X(s)=\EE e^{\ii sX}\), then
\begin{equation}\label{eq:first-channel-laplace-resolvent}
 \frac{1+c_1(t)}{2t}
 =\EE\frac1{2t-\ii X}
 =\int_0^\infty e^{-2ts}\varphi_X(s)\,\dd s.
\end{equation}
Thus the full scale trajectory identifies the centred law by the classical
uniqueness of the Laplace and characteristic transforms.  It identifies the
untranslated law \(\nu\) only when \(\bar\gamma\) is supplied separately.

Moreover, if \(n\ge1\) and \(\mu_{2j}=\EE X^{2j}<\infty\) for
\(1\le j<n\), then
\begin{equation}\label{eq:first-channel-even-moment-remainder}
\begin{aligned}
 &(-1)^{n-1}4^nt^{2n}
 \left[-\Re c_1(t)-\sum_{j=1}^{n-1}
 (-1)^{j-1}\frac{\mu_{2j}}{4^jt^{2j}}\right]\\
 &\hspace{32mm}
 =\EE\!\left[\frac{X^{2n}}{1+X^2/(4t^2)}\right]
 \uparrow\EE|X|^{2n}
 \quad(t\to\infty),
\end{aligned}
\end{equation}
where the limit is permitted to be \(+\infty\).
\end{theorem}

\begin{proof}
Fix a deterministic translation \(\varepsilon=X/t\).
The pole calculation in
Lemma~\ref{lem:second-mode-cauchy-translate-transform} gives
\[
 R_\varepsilon(y)
 =1+\sum_{k\ge1}\bigl(\overline{a_t(X)}^{\,k}z^k
 +a_t(X)^kz^{-k}\bigr).
\]
Indeed, the interior pole in that calculation is
\(-\varepsilon/(\varepsilon+2\ii)=\ii X/(2t-\ii X)\), while the inverse of
the exterior pole is its complex conjugate.  For finite \(X\),
\(|a_t(X)|<1\).  On a compact subset of
\(\TT\setminus\{-1\}\), the denominators in the two geometric remainders
are bounded away from zero uniformly over the compactified parameter line
\(X\in\RR\cup\{\infty\}\); the only boundary value of \(a_t(X)\) is
\(-1\).  Dominated convergence therefore permits expectation of the finite
partial sums and passage to the limit, proving
\eqref{eq:first-channel-exact-fourier-series} and the asserted local
uniformity.

Pointwise differentiation gives
\[
 a_t(X)^{k+1}=-a_t(X)^k-\frac{t}{k}
 \frac{\dd}{\dd t}a_t(X)^k.
\]
The derivative is integrable locally uniformly in \(t\), because the same
identity bounds its absolute value by \(2k/t\).  Differentiation under the
expectation proves \eqref{eq:first-channel-differential-closure}.

Also \(1+a_t(X)=2t/(2t-\ii X)\).  The identity
\[
 \frac1{2t-\ii X}=\int_0^\infty e^{-(2t-\ii X)s}\,\dd s
\]
and Fubini's theorem prove \eqref{eq:first-channel-laplace-resolvent}.

If \(\EE X^2<\infty\), Jensen's inequality and the geometric series give
\[
 \sum_{k\ge1}|c_k(t)|^2
 \le\EE\sum_{k\ge1}|a_t(X)|^{2k}
 =\frac{\EE X^2}{4t^2}<\infty.
\]
Thus the Fourier series converges in \(L^2(\omega)\), where it agrees with
the quotient already identified pointwise.  Parseval's identity gives
\eqref{eq:first-channel-chi-square-parseval}.

Finally,
\begin{equation}\label{eq:first-channel-real-part}
 -\Re c_1(t)=\EE\frac{X^2}{4t^2+X^2}.
\end{equation}
For \(q\ge0\), finite geometric division yields
\[
 \frac q{1+q}=\sum_{j=1}^{n-1}(-1)^{j-1}q^j
 +(-1)^{n-1}\frac{q^n}{1+q}.
\]
Apply this identity with \(q=X^2/(4t^2)\), take expectations, and multiply
by \(4^nt^{2n}\).  The right side of
\eqref{eq:first-channel-even-moment-remainder} increases pointwise with
\(t\), so monotone convergence proves the final assertion.
\end{proof}

\subsection{Moment-matched comparisons and Gaussian compression}

For a probability law \(\rho\) with a finite \(j\)th absolute moment, write
\(m_j(\rho)=\int x^j\rho(\dd x)\).  The comparison theorem below does not
require either law to be centred.

\begin{theorem}[First unmatched moment for Poisson relative entropy]
\label{thm:moment-matched-poisson-kl}
Let \(r\ge2\), and let \(\nu,\eta\) be probability laws on \(\RR\) with
finite \(r\)th absolute moments.  Suppose
\[
 m_j(\nu)=m_j(\eta)\quad(0\le j<r),
 \qquad \Delta_r:=m_r(\nu)-m_r(\eta)\ne0.
\]
Then
\begin{equation}\label{eq:moment-matched-poisson-kl}
 D_{\rm KL}(P_t*\nu\|P_t*\eta)
 =C_r\Delta_r^2t^{-2r}+o(t^{-2r}),
 \qquad
 C_r=4^{-r}\binom{2r-2}{r-1}.
\end{equation}
\end{theorem}

\begin{proof}
Divide both smoothed densities by \(P_t(ty)\), and denote the resulting
quotients by \(q_t^\nu\) and \(q_t^\eta\).  The global two-region Taylor
remainder \eqref{eq:one-dimensional-two-region-taylor-remainder}, applied
at order \(r\), gives
\begin{equation}\label{eq:moment-matched-quotient-difference}
 q_t^\nu-q_t^\eta
 =\Delta_r u_r t^{-r}+o(t^{-r})
 \quad\text{in }L^\infty(\omega)\cap L^1(\omega).
\end{equation}
The same argument at order one shows \(q_t^\eta\to1\) uniformly.  Since the
two quotient differences have zero \(\omega\)-mean, the exact Bregman form is
\[
 D_{\rm KL}(P_t*\nu\|P_t*\eta)
 =\int q_t^\eta
 \Phi\!\left(\frac{q_t^\nu-q_t^\eta}{q_t^\eta}\right)\dd\omega.
\]
Uniform Taylor expansion of \(\Phi(s)=s^2/2+O(s^3)\), together with
\eqref{eq:moment-matched-quotient-difference}, yields
\[
 D_{\rm KL}(P_t*\nu\|P_t*\eta)
 =\frac{\Delta_r^2}{2}t^{-2r}\int u_r^2\dd\omega+o(t^{-2r}).
\]
By the Laurent coefficient formula
\eqref{eq:mode-laurent-coefficients} and Parseval,
\[
 \frac12\int u_r^2\dd\omega
 =4^{-r}\sum_{j=0}^{r-1}\binom{r-1}{j}^2
 =4^{-r}\binom{2r-2}{r-1},
\]
where the last equality is Vandermonde's identity.
\end{proof}

We now state the Gaussian quadrature consequence with the nondegeneracy
hypothesis made explicit.  This qualification is necessary: an \(n\)-node
quadrature rule is not defined for an arbitrary law supported on fewer than
\(n\) points.

\begin{theorem}[Gaussian-comparison entropy threshold]
\label{thm:gauss-poisson-kl-threshold}
Let \(n\ge1\), and let \(\nu\) be a probability law on \(\RR\) satisfying
\[
 \EE|X|^{2n-1}<\infty,
 \qquad |\operatorname{supp}\nu|\ge n.
\]
Let \(\pi_n\) be the monic degree-\(n\) polynomial orthogonal to
\(1,x,\ldots,x^{n-1}\) with respect to \(\nu\), and let \(G_n\nu\) be its
\(n\)-node Gaussian quadrature law.  Thus
\begin{equation}\label{eq:gauss-moment-matching}
 m_j(G_n\nu)=m_j(\nu),\qquad 0\le j\le2n-1.
\end{equation}
Then
\begin{equation}\label{eq:gauss-poisson-finite-limit}
 \EE|X|^{2n}<\infty
 \quad\Longrightarrow\quad
 \lim_{t\to\infty}t^{4n}
 D_{\rm KL}(P_t*\nu\|P_t*G_n\nu)
 =4^{-2n}\binom{4n-2}{2n-1}
 \|\pi_n\|_{L^2(\nu)}^4,
\end{equation}
whereas
\begin{equation}\label{eq:gauss-poisson-infinite-limit}
 \EE|X|^{2n}=\infty
 \quad\Longrightarrow\quad
 t^{4n}D_{\rm KL}(P_t*\nu\|P_t*G_n\nu)\longrightarrow+\infty.
\end{equation}
Consequently,
\begin{equation}\label{eq:gauss-poisson-moment-equivalence}
 \EE|X|^{2n}<\infty
 \quad\Longleftrightarrow\quad
 \limsup_{t\to\infty}t^{4n}
 D_{\rm KL}(P_t*\nu\|P_t*G_n\nu)<\infty.
\end{equation}
\end{theorem}

\begin{proof}
The existence, positivity, and exactness
\eqref{eq:gauss-moment-matching} of the Gaussian rule are classical; see
\cite{GolubWelsch1969GaussQuadrature}.  If the \(2n\)th moment is finite,
then, because \(\pi_n\) vanishes at every quadrature node and
\(\pi_n^2-x^{2n}\) has degree at most \(2n-1\),
\begin{equation}\label{eq:gauss-missing-moment-norm}
 m_{2n}(\nu)-m_{2n}(G_n\nu)
 =\int\pi_n^2\dd\nu=\|\pi_n\|_{L^2(\nu)}^2.
\end{equation}
Theorem~\ref{thm:moment-matched-poisson-kl}, with \(r=2n\), now proves
\eqref{eq:gauss-poisson-finite-limit}.  If the norm in
\eqref{eq:gauss-missing-moment-norm} vanishes, then \(\nu\) is itself
supported on the Gaussian nodes and both sides of
\eqref{eq:gauss-poisson-finite-limit} are zero.

Suppose instead that \(\EE|X|^{2n}=\infty\).  Apply
\eqref{eq:first-channel-even-moment-remainder} to \(\nu\) and to the
finite-support law \(G_n\nu\).  Their even moments through order \(2n-2\)
agree, and therefore
\[
 t^{2n}\left|\Re\{c_1^\nu(t)-c_1^{G_n\nu}(t)\}\right|
 \longrightarrow+\infty.
\]
The real part of \(z\) is a bounded test of norm one on the Cayley circle.
Pinsker's inequality and the Fourier coefficient identity therefore give
\[
 2D_{\rm KL}(P_t*\nu\|P_t*G_n\nu)
 \ge\left|\Re\{c_1^\nu(t)-c_1^{G_n\nu}(t)\}\right|^2.
\]
Multiplication by \(t^{4n}\) proves
\eqref{eq:gauss-poisson-infinite-limit}, and the equivalence follows.
\end{proof}

\begin{remark}[Classical optimal degree]
The assertion that an \(n\)-node positive quadrature rule has maximal exact
degree \(2n-1\), and that Gaussian quadrature attains this degree, is a
classical quadrature theorem, not a new result of this article.  The node
polynomial gives the elementary obstruction at degree \(2n\); see
\cite{GolubWelsch1969GaussQuadrature}.
\end{remark}

\begin{corollary}[Two-node fourth-moment coefficient]
\label{cor:two-node-gauss-fourth-moment}
Assume
\[
 \EE X=0,\qquad \EE X^2=1,\qquad \EE|X|^3<\infty,
\]
and put
\[
 \beta_3=\EE X^3,\qquad
 \pi_2(x)=x^2-\beta_3x-1.
\]
Let \(G_2\nu\) be the two-node Gaussian quadrature law.  If
\(\EE X^4<\infty\), then, with
\[
 \kappa=\|\pi_2\|_{L^2(\nu)}^2
 =\EE X^4-1-\beta_3^2,
\]
one has
\begin{equation}\label{eq:two-node-fourth-moment-coefficient}
 \lim_{t\to\infty}t^8D_{\rm KL}(P_t*\nu\|P_t*G_2\nu)
 =\frac5{64}\kappa^2.
\end{equation}
If \(\EE X^4=\infty\), the left side of
\eqref{eq:two-node-fourth-moment-coefficient} tends to \(+\infty\).  For a
symmetric law,
\(G_2\nu=(\delta_{-1}+\delta_1)/2\) and
\(\kappa=\EE X^4-1\).
\end{corollary}

\begin{proof}
The polynomial \(\pi_2\) is monic and orthogonal to \(1,x\).  Its squared
norm is the identity already recorded in
Lemma~\ref{lem:appendix-normalized-fourth-moment-defect-algebra}.  Apply
Theorem~\ref{thm:gauss-poisson-kl-threshold} and use
\(4^{-4}\binom63=5/64\).  In the symmetric case, exactness through degree
three and unit variance force the nodes and weights to be
\(\{-1,1\}\) and \(\{1/2,1/2\}\).
\end{proof}

\begin{theorem}[Square boundary at a missing Gaussian moment]
\label{thm:gauss-regular-variation-square-law}
Let \(n\ge2\), put \(p=2n\), and let \(\nu\) satisfy the hypotheses of
Theorem~\ref{thm:gauss-poisson-kl-threshold}.  Assume, for an eventually
positive slowly varying function \(L\) and constants \(c_+,c_-\ge0\) with
\(c_++c_->0\),
\[
 \mathbb P(X>x)\sim c_+x^{-p}L(x),\qquad
 \mathbb P(X<-x)\sim c_-x^{-p}L(x),
\]
and assume
\[
 \ell_L(t)=\int_{x_0}^t\frac{L(s)}s\dd s\longrightarrow\infty.
\]
Then, with
\(M_p(t)=\EE[|X|^p\mathbf1_{\{|X|\le t\}}]\),
\begin{equation}\label{eq:gauss-regular-variation-square-law}
 D_{\rm KL}(P_t*\nu\|P_t*G_n\nu)
 \sim C_p t^{-2p}M_p(t)^2,
 \qquad C_p=4^{-p}\binom{2p-2}{p-1}.
\end{equation}
In particular,
\[
 M_p(t)\sim p(c_++c_-)\ell_L(t),
\]
so the boundary normalization is the square of the integrated slowly
varying factor and the leading constant is positive.
\end{theorem}

\begin{proof}
The quotient estimate proved for
Theorem~\ref{cor:main-body-conditional-missing-top-moment-tail}, before the
entropy transfer is applied, gives
\[
 q_t^\nu
 =\sum_{j=0}^{p-1}m_j(\nu)u_jt^{-j}
 +u_pt^{-p}M_p(t)+e_t,
 \qquad
 \|e_t\|_\infty+\|e_t\|_{L^1(\omega)}
 =o(t^{-p}M_p(t)).
\]
The finite-support law \(G_n\nu\) has an ordinary expansion through order
\(p\), and its moments through order \(p-1\) agree with those of \(\nu\).
Since \(M_p(t)\to\infty\), subtraction yields
\begin{equation}\label{eq:gauss-boundary-quotient-difference}
 q_t^\nu-q_t^{G_n\nu}
 =u_pt^{-p}M_p(t)+o_{L^\infty\cap L^1(\omega)}
 (t^{-p}M_p(t)).
\end{equation}
Interpolation also gives the same relative error in \(L^2(\omega)\).
Furthermore \(q_t^{G_n\nu}\to1\) uniformly.  Applying the Bregman expansion
from the proof of Theorem~\ref{thm:moment-matched-poisson-kl} to
\eqref{eq:gauss-boundary-quotient-difference} gives
\[
 D_{\rm KL}(P_t*\nu\|P_t*G_n\nu)
 =\frac12t^{-2p}M_p(t)^2\int u_p^2\dd\omega
 +o(t^{-2p}M_p(t)^2).
\]
The Parseval calculation in that theorem identifies the constant as \(C_p\).
The asymptotic for \(M_p\) is Karamata's boundary integral theorem, exactly
as in Theorem~\ref{cor:main-body-conditional-missing-top-moment-tail}.
\end{proof}

\subsection{The multidimensional \texorpdfstring{\(L^2\)}{L2} and entropy
thresholds}

Write \(a=(d+1)/2\) and
\[
 F_y(z)=\left(\frac{1+|y|^2}{1+|y-z|^2}\right)^a.
\]
Let \(T_2(y,z)\) denote the Taylor polynomial of \(F_y(z)\) at \(z=0\)
through homogeneous degree two.

\begin{lemma}[Large translations in the Cauchy quotient space]
\label{lem:multidim-large-translation-l2-growth}
For every \(d\ge1\),
\begin{equation}\label{eq:translated-quotient-l2-asymptotic}
 \|F_\cdot(z)\|_{L^2(\Omega_d)}^2
 \sim2^{-d}|z|^{d+1}\qquad(|z|\to\infty).
\end{equation}
If
\[
 q_d:=\max\left\{2,\frac{d+1}{2}\right\},
\]
then a constant \(C_d\) exists such that
\begin{equation}\label{eq:second-order-l2-global-remainder}
 \|F_\cdot(z)-T_2(\cdot,z)\|_{L^2(\Omega_d)}
 \le C_d
 \begin{cases}|z|^3,&|z|\le1,\\ |z|^{q_d},&|z|>1.
 \end{cases}
\end{equation}
\end{lemma}

\begin{proof}
With \(c_d=\Gamma((d+1)/2)/\pi^{(d+1)/2}\) and \(y=z+u\),
\[
 \|F_\cdot(z)\|_2^2
 =c_d\int_{\RR^d}
 \frac{(1+|z+u|^2)^a}{(1+|u|^2)^{d+1}}\dd u.
\]
After division by \(|z|^{d+1}\), dominated convergence applies and gives
\[
 c_d\int_{\RR^d}(1+|u|^2)^{-(d+1)}\dd u
 =\frac{\Gamma((d+1)/2)\Gamma(d/2+1)}
 {\sqrt\pi\Gamma(d+1)}=2^{-d}.
\]
For \(|z|\le1\), differentiation under the \(L^2(\Omega_d)\) norm through
third order is uniform on the unit ball, and Taylor's formula gives the
first line of \eqref{eq:second-order-l2-global-remainder}.  For \(|z|>1\),
\eqref{eq:translated-quotient-l2-asymptotic}, enlarged to a global upper
bound, and the fact that \(T_2\) is a degree-two polynomial in \(z\) with
\(L^2(\Omega_d)\) coefficients give
\[
 \|F_\cdot(z)-T_2(\cdot,z)\|_2
 \le C_d(|z|^{(d+1)/2}+|z|^2),
\]
which is the second line.
\end{proof}

\begin{theorem}[Sharp absolute-moment condition for the \(L^2\) quotient]
\label{thm:multidim-l2-moment-threshold}
Let \(\nu\) be a centred probability law on \(\RR^d\), with covariance
matrix \(\Sigma\).  If
\begin{equation}\label{eq:multidim-qd-moment}
 \EE|X|^{q_d}<\infty,
 \qquad q_d=\max\left\{2,\frac{d+1}{2}\right\},
\end{equation}
then
\begin{equation}\label{eq:multidim-l2-second-order-convergence}
 t^2\delta_t^{(d)}\longrightarrow b_\Sigma
 \quad\text{in }L^2(\Omega_d).
\end{equation}
For \(d\ge4\), this exponent is sharp among hypotheses stated solely as a
finite absolute moment: for every \(2\le p<q_d\), there is a centred law
with positive-definite covariance and finite \(p\)th absolute moment for
which \eqref{eq:multidim-l2-second-order-convergence} fails.
\end{theorem}

\begin{proof}
After centering removes the degree-one term and covariance averages the
degree-two term, Minkowski's inequality gives
\[
 \|t^2\delta_t^{(d)}-b_\Sigma\|_2
 \le\EE\left[t^2
 \|F_\cdot(X/t)-T_2(\cdot,X/t)\|_2\right].
\]
The integrand converges pointwise to zero.  On \(|X|\le t\), the first line
of \eqref{eq:second-order-l2-global-remainder} bounds it by \(C|X|^2\);
on \(|X|>t\), the second line bounds it by
\(C|X|^{q_d}t^{2-q_d}\le C|X|^{q_d}\).  Dominated convergence proves
\eqref{eq:multidim-l2-second-order-convergence}.

For sharpness, let \(a=(d+1)/2>2\), choose \(p<a\), and take a symmetric
spike law with masses \(w_j/2\) at \(\pm R_je_1\), where
\(w_j=\alpha_jR_j^{-p}\), \(\sum_j\alpha_j<\infty\), and \(R_j\) grows
recursively.  Add a sufficiently small bounded symmetric component spanning
\(\RR^d\).  The resulting law has positive-definite finite covariance and
finite \(p\)th moment.  Choose
\[
 0<s<\min\left\{\frac{a-p}{a-2},\,1-\frac{p}{d+1}\right\}
\]
and put \(t_j=R_j^s\).
On a fixed ball about \(R_je_1/t_j\), the \(j\)th translated quotient has
\(L^2(\Omega_d)\) size comparable to
\((R_j/t_j)^a\), by the localized proof of
\eqref{eq:translated-quotient-l2-asymptotic}.  Consequently its contribution
to \(t_j^2\delta_{t_j}^{(d)}\) has size at least
\[
 c\alpha_jR_j^{a-p-s(a-2)}.
\]
The first restriction on \(s\) makes the displayed exponent positive; the
second gives \(w_j(R_j/t_j)^{d+1}\to\infty\), so the localized positive spike
also dominates the constant background.  Choosing \(R_j\) fast enough makes
the displayed quantity diverge and dominates the bounded quadratic mode and
all previously chosen terms.  Hence the \(L^2\)
convergence fails.
\end{proof}

\begin{theorem}[Low-dimensional finite-covariance entropy threshold]
\label{thm:low-dimensional-finite-covariance-threshold}
Let \(1\le d\le3\), and let \(\nu\) be a probability law on \(\RR^d\) with
finite mean.  With
\[
 H_d(t)=D_{\rm KL}\!\left(P_t^{(d)}*\nu\,\middle\|\,
 P_t^{(d)}(\mathord\cdot-\bar\gamma)\right),
\]
one has
\begin{equation}\label{eq:low-dimensional-covariance-iff}
 \EE|X|^2<\infty
 \quad\Longleftrightarrow\quad
 \limsup_{t\to\infty}t^4H_d(t)<\infty.
\end{equation}
If the covariance is finite, then
\begin{equation}\label{eq:low-dimensional-covariance-limit}
 H_d(t)=\mathcal Q_d(\Sigma)t^{-4}+o(t^{-4}).
\end{equation}
If it is infinite, then \(t^4H_d(t)\to+\infty\) in the extended sense.
\end{theorem}

\begin{proof}
For \(d\le3\), \(q_d=2\).  Thus Theorem
\ref{thm:multidim-l2-moment-threshold} and the \(L^2\) alternative in
Theorem~\ref{thm:second-order-quotient-kl-coefficient} prove
\eqref{eq:low-dimensional-covariance-limit}.  The entropy representation is
valid in the extended sense by change of variables; for large \(t\), the
\(L^2\) convergence makes \(\Phi(\delta_t^{(d)})\) integrable because
\(\Phi(s)\le C s^2\) on \([-1,\infty)\).

Conversely, if \(\EE|X|^2=\infty\), some unit vector \(e\) satisfies
\(\EE(e\cdot X)^2=\infty\).  Projection onto \(e\) sends the radial
\(d\)-dimensional Poisson law to the one-dimensional Cauchy law at the same
scale.  The data-processing inequality and
Theorem~\ref{thm:sharp-leading-kl-variance-threshold} therefore give
\[
 H_d(t)\ge
 D_{\rm KL}(P_t*\nu_e\|P_t(\mathord\cdot-\bar\gamma_e)),
\]
whose product with \(t^4\) tends to \(+\infty\).
\end{proof}

\begin{proposition}[Finite covariance does not control high-dimensional KL]
\label{prop:high-dimensional-finite-covariance-obstruction}
For every \(d\ge4\), there exists a centred discrete probability law on
\(\RR^d\) with positive-definite finite covariance such that
\begin{equation}\label{eq:high-dimensional-finite-covariance-obstruction}
 \limsup_{t\to\infty}t^4H_d(t)=+\infty.
\end{equation}
\end{proposition}

\begin{proof}
Choose positive \(\alpha_j\) with \(\sum_j\alpha_j<\infty\), radii
\(R_j\uparrow\infty\), and weights
\(w_j=\alpha_jR_j^{-2}\).  After rescaling the sequence if necessary,
\(W=\sum_jw_j<1\).  Begin with
\[
 (1-W)\delta_0+\sum_{j\ge1}\frac{w_j}{2}
 (\delta_{R_je_1}+\delta_{-R_je_1}),
\]
and transfer an arbitrarily small amount of the mass at zero to a bounded
centred symmetric law spanning \(\RR^d\).  The covariance is then positive
definite and finite because \(\sum_jw_jR_j^2=\sum_j\alpha_j<\infty\).

Fix a sufficiently small constant \(\kappa>0\), and set
\[
 t_j=\kappa R_jw_j^{1/(d+1)}.
\]
The radii may be chosen recursively so that \(t_j\to\infty\).  Let \(A_j\)
be the ball of radius \(t_j\) centred at \(R_je_1\).  The \(j\)th mixture
component gives
\[
 (P_{t_j}^{(d)}*\nu)(A_j)\ge c_dw_j,
\]
whereas, since \(t_j/R_j\to0\), the centred reference satisfies
\[
 P_{t_j}^{(d)}(A_j)
 \le C_d(t_j/R_j)^{d+1}
 =C_d\kappa^{d+1}w_j.
\]
Choose \(\kappa\) so that the second constant is strictly smaller than the
first.  Data processing to the partition \(\{A_j,A_j^c\}\) and the
elementary Bernoulli relative-entropy bound then give
\(H_d(t_j)\ge c'_dw_j\).  Hence
\[
 t_j^4H_d(t_j)
 \ge c''_dR_j^4w_j^{1+4/(d+1)}
 =c''_d\alpha_j^{1+4/(d+1)}
 R_j^{2(d-3)/(d+1)}.
\]
The exponent of \(R_j\) is positive exactly when \(d>3\).  Choosing the
radii recursively makes the final expression tend to infinity.
\end{proof}

\begin{remark}[Open KL moment interface]
\label{rem:optimal-kl-moment-interface}
The preceding localized argument also gives a rigorous obstruction below
\[
 p_{\rm KL}(d)=\frac{4(d+1)}{d+5}.
\]
Indeed, for every \(d\ge4\) and
\(2\le p<p_{\rm KL}(d)\), weights
\(w_j=\alpha_jR_j^{-p}\) produce a finite \(p\)th moment while the lower
bound scales as
\[
 R_j^{\,4-p(d+5)/(d+1)}
\]
up to a positive power of \(\alpha_j\); sufficiently rapid radii make
\(t_j^4H_d(t_j)\) unbounded.  What remains open is the load-bearing
sufficiency statement at and above this exponent: it is not proved here
that
\[
 \EE|X|^{\max\{2,p_{\rm KL}(d)\}}<\infty
 \quad\Longrightarrow\quad
 H_d(t)=\mathcal Q_d(\Sigma)t^{-4}+o(t^{-4}).
\]
Thus \(p_{\rm KL}(d)\) is presently a sharp candidate forced by the spike
obstruction, not an established KL threshold.
\end{remark}
